\section{Block Accumulate Minimum Distance Proof}\label{sec:BAProofApd}



\if\submission1
\else
\paragraph{Proof overview.}
Fix a nonzero input $\xv\in\{0,1\}^k$ and track Hamming weights through
$\G=\B\pi_1\A\pi_2\A$:
\[
u:=\hw{\xv\B},\qquad
h_1:=\hw{\xv\B\pi_1\A},\qquad
h:=\hw{\xv\B\pi_1\A\pi_2\A}.
\]
We upper bound the expected number of nonzero codewords with final weight
$h\le \varepsilon n$ and then apply Markov's inequality.
\begin{enumerate}
\item \emph{Outer block-diagonal distance.}
  By \eqref{eq:umin_def}, any weight-$w$ input activates at least
  $\lceil w/\sigma\rceil$ blocks, and each active block contributes at least $d'$
  output weight. Hence
  $u=\hw{\xv\B}\ge u_{\min}(w)=d'\lceil w/\sigma\rceil$.

\item  \emph{Random interleavers give uniformity.}
  Conditioned on $\hw{\xv\B}=u$, the vector after $\pi_1$ is uniform among all
  length-$n$ vectors of weight $u$; similarly, conditioned on $h_1$, the vector
  after $\pi_2$ is uniform among all weight-$h_1$ vectors. Thus both accumulator
  stages are analyzed via the exact input-output enumerator $\enum_{u,h}$ of the accumulator.

\item  \emph{Accumulators suppress very small outputs.}
  Lemma~\ref{lem:small_eps} shows that for $h_1\le 2\varepsilon n$,
  \[
  \Pr[\hw{z\A}=h_1\mid \hw{z}=u]\;\le\;\left(4e^2\frac{h_1}{n}\right)^{u/2},
  \]
  and the RHS is monotone decreasing in $u$ when $\varepsilon$ is small. This
  lets us eliminate the (potentially complicated) distribution of $u$ given
  $\xv$ and upper bound by the worst case allowed by \eqref{eq:umin_def}.
  Applying the same estimate to the second accumulator yields a strong tail bound
  on $\Pr[h\le \varepsilon n\mid h_1]$.

\item \emph{Support constraints truncate the weight ranges.}
  Since $\enum_{a,b}=0$ unless $\lceil a/2\rceil\le b$, any final weight
  $h\le \varepsilon n$ forces $h_1\le 2\varepsilon n$ and in turn forces
  $u\le 2h_1$. Combined with $u\ge u_{\min}(w)$, this restricts the relevant
  inputs to a low-weight regime.

\item \emph{Re-index by the number of active blocks.}
  The key counting step is to sum over inputs by the number $b$ of active
  $\sigma$-bit blocks rather than by total Hamming weight $w$.
  For inputs with exactly $b$ active blocks we have the sharp bound
  $u\ge d'b$, and the number of such inputs is at most
  $\binom{k/\sigma}{b}(2^\sigma-1)^b$.
  This block-based counting aligns with the structure of $\B$ and yields a
  geometric decay in $b$ once $d'>4$, implying that the expected number of
  codewords with $h\le \varepsilon n$ tends to zero.
\end{enumerate}

\fi




\begin{assumption}[Outer block-diagonal distance]
Assume $\B$ is block-diagonal with input-block size $\sigma$ and that each activated block produces output Hamming weight at least some constant $d'>4$. Then for any $\xv$ with $w:=\hw{\xv}$,
\begin{equation}
\label{eq:umin_def}
 u := \hw{\xv\B} \;\ge\; u_{\min}(w) := d' \Big\lceil \frac{w}{\sigma}\Big\rceil.
\end{equation}
\end{assumption}


\paragraph{Accumulator input-output enumerator.}
Using the standard accumulator from the preliminaries, we recall the (exact) input--output enumerator:
\begin{equation}
\label{eq:acc_exact}
\enum^\A_{u,h}
=
\binom{n-h}{\lfloor u/2\rfloor}\binom{h-1}{\lceil u/2\rceil-1}.
\end{equation}
We will omit the superscript \A in this section. In particular, for uniform $z\in\{0,1\}^n$ conditioned on $\hw{z}=u$,
\begin{equation}
\label{eq:acc_prob_exact}
\Pr[\hw{z\A}=h \mid \hw{z}=u]
=
\frac{\enum_{u,h}}{\binom{n}{u}},
\end{equation}
and $\enum_{u,h}=0$ unless $\lceil u/2\rceil \le h$.



\begin{lemma}[Small-output contraction for the accumulator]
\label{lem:small_eps}
Let $\A$ be the length-$n$ accumulator with enumerator \eqref{eq:acc_exact}.
There exists a constant $\varepsilon_0>0$ such that for all
$0<\varepsilon\le \varepsilon_0$, all sufficiently large $n$, all
$1\le h_1\le 2\varepsilon n$, and all $u\ge 1$,
\begin{equation}
\label{eq:small_eps_i}
\Pr\big[\hw{z\A}=h_1\mid \hw{z}=u\big]
=\frac{\enum_{u,h_1}}{\binom{n}{u}}
\le \left(4e^2\frac{h_1}{n}\right)^{u/2}.
\end{equation}
Moreover, if $8e^2\varepsilon<1$, then the RHS of \eqref{eq:small_eps_i} is
decreasing in $u$.

In addition, for each fixed $h_1\le 2\varepsilon n$,
\begin{equation}
\label{eq:small_eps_ii}
\sum_{h=1}^{\lfloor \varepsilon n\rfloor}
\Pr\big[\hw{z\A}=h\mid \hw{z}=h_1\big]
\le (\varepsilon n)\left(4e^2\varepsilon\right)^{h_1/2}.
\end{equation}
\end{lemma}

\begin{proof}
We prove \eqref{eq:small_eps_i} first. Let $s=\lfloor u/2\rfloor$ and
$t=\lceil u/2\rceil$, so $s+t=u$ and $s,t\ge u/2-1/2$. Using
\eqref{eq:acc_exact},
$
\enum_{u,h_1}
=
\binom{n-h_1}{s}\binom{h_1-1}{t-1}
\le
\binom{n}{s}\binom{h_1}{t}.
$
Apply $\binom{N}{r}\le (eN/r)^r$ to both factors:
$
\enum_{u,h_1}
\le
\left(\frac{en}{s}\right)^s\left(\frac{eh_1}{t}\right)^t.
$
Also use $\binom{n}{u}\ge (n/u)^u$. Then
$
\frac{\enum_{u,h_1}}{\binom{n}{u}}
\le
\left(\frac{en}{s}\right)^s\left(\frac{eh_1}{t}\right)^t\left(\frac{u}{n}\right)^u
=
\left(e\frac{u}{s}\right)^s\left(e\frac{u}{t}\right)^t\left(\frac{h_1}{n}\right)^t.
$
If $u=1$, then $\enum_{1,h_1}=1$ and $\binom{n}{1}=n$, so
$\Pr[\hw{z\A}=h_1\mid \hw{z}=1]=1/n\le (4e^2(h_1/n))^{1/2}$ for all sufficiently large $n$.
Assume henceforth that $u\ge 2$, so $s,t\ge 1$.
If $u$ is even then $s=t=u/2$ and $(u/s)^s(u/t)^t=2^u$.
If $u$ is odd, write $u=2m+1$ so $s=m$ and $t=m+1$, and observe that
\[
\left(\frac{u}{s}\right)^s\left(\frac{u}{t}\right)^t
=
2^u\left(1+\frac{1}{2m}\right)^m\left(1-\frac{1}{2m+2}\right)^{m+1}
\le 2^u,
\]
where the last inequality follows from $\log(1+x)\le x$ and $\log(1-x)\le -x-x^2/2$.
Thus $\left(e\frac{u}{s}\right)^s\left(e\frac{u}{t}\right)^t \le (2e)^u$, and hence
\[
\frac{\enum_{u,h_1}}{\binom{n}{u}}
\le
(2e)^{u}\left(\frac{h_1}{n}\right)^{t}
\le
(2e)^{u}\left(\frac{h_1}{n}\right)^{u/2}.
\]


Finally, $(2e)^u=(4e^2)^{u/2}$ gives \eqref{eq:small_eps_i}. If
$h_1\le 2\varepsilon n$ and $8e^2\varepsilon<1$, then
$4e^2(h_1/n)\le 8e^2\varepsilon<1$, so the RHS of \eqref{eq:small_eps_i} is
decreasing in $u$.

For \eqref{eq:small_eps_ii}, apply \eqref{eq:small_eps_i} with $u\leftarrow h_1$
to get, for each $h$,
$
\Pr[\hw{z\A}=h\mid \hw{z}=h_1]
\le
\left(4e^2\frac{h}{n}\right)^{h_1/2}.
$
Summing over $1\le h\le \varepsilon n$ and using $h/n\le \varepsilon$ yields
\[
\sum_{h=1}^{\lfloor \varepsilon n\rfloor}
\Pr[\hw{z\A}=h\mid \hw{z}=h_1]
\le
(\varepsilon n)\left(4e^2\varepsilon\right)^{h_1/2}.
\]
\end{proof}


\begin{theorem}[Linear minimum distance for BA-$3$]
\label{thm:ba_linear_distance}
Consider the Block-Accumulate code $\G:=\B\pi_1\A\pi_2\A$ of length $n$, where $\pi_1,\pi_2$ are independent uniform interleavers on $[n]$ and $\A$ is the standard length-$n$ accumulator.
Assume the outer block-diagonal distance lower bound \eqref{eq:umin_def} holds.
If the outer block distance satisfies $d'>4$, then for some sufficiently small
constant $\varepsilon=\varepsilon(d',\sigma,R)>0$,
\[
\Pr[d_{\min}(\G)\ge \varepsilon n]\xrightarrow[n\to\infty]{}1.
\]
\end{theorem}
\begin{remark}
  We write $R:=k/n$ for the (constant) rate of the code. Here $\varepsilon$ depends only on $d', \sigma$ and the rate $R$, and not on $n$.
\end{remark}
\begin{proof}

%------------------------------------------------------------
\noindent\textbf{Step 1: Expected number of low-weight codewords.}

Fix a constant $0<\varepsilon\le \varepsilon_0$ sufficiently small so that
$8e^2\varepsilon<1$, $(4e^2\varepsilon)^{1/2}e<1$, and $\rho<1$ (as defined later).  Let $E_{\le \varepsilon n}$ denote the expected number of \emph{nonzero}
codewords of weight at most $\varepsilon n$.
By linearity of expectation,
\begin{equation}
\label{eq:Elow_start}
E_{\le \varepsilon n}
=
\sum_{w=1}^{k} \binom{k}{w}\;
\Pr\big[\hw{\xv\G}\le \varepsilon n\big],
\end{equation}
where $\xv$ is uniform over $\{x\in\{0,1\}^k:\hw{x}=w\}$.

Introduce the intermediate weights
\[
 u=\hw{\xv\B},\qquad
h_1=\hw{\xv\B\pi_1\A},\qquad
h=\hw{\xv\B\pi_1\A\pi_2\A}=\hw{\xv\G}=\hw{\yv}.
\]
Conditioning on $(u,h_1,h)$,
\begin{equation}
\label{eq:Elow_expand}
E_{\le \varepsilon n}
=
\sum_{w=1}^{k} \binom{k}{w}
\sum_{u=0}^{n}\sum_{h_1=0}^{n}\sum_{h=1}^{\lfloor \varepsilon n\rfloor}
\Pr[u \mid w]\;
\Pr[h_1 \mid u]\;
\Pr[h \mid h_1].
\end{equation}

\noindent\textbf{Uniformity under random interleavers.}
Conditioned on $\hw{\xv\B}=u$, the vector $\xv\B\pi_1$ is uniform over all
weight-$u$ vectors in $\{0,1\}^n$; similarly, conditioned on $h_1$, the vector
$\xv\B\pi_1\A\pi_2$ is uniform over all weight-$h_1$ vectors in $\{0,1\}^n$.
Thus \eqref{eq:acc_prob_exact} applies to both accumulator stages:
\begin{equation}
\label{eq:stage_probs}
\Pr[h_1\mid u]=\frac{\enum_{u,h_1}}{\binom{n}{u}},
\qquad
\Pr[h\mid h_1]=\frac{\enum_{h_1,h}}{\binom{n}{h_1}}.
\end{equation}

%------------------------------------------------------------
\noindent\textbf{Step 2: Remove the distribution of $u$ using monotonicity.}

We only know $u\ge u_{\min}(w)$ from \eqref{eq:umin_def}; the exact distribution
$\Pr[u\mid w]$ may be complicated. We eliminate it using a one-sided bound.

For $h_1\le 2\varepsilon n$, Lemma~\ref{lem:small_eps}(i) gives
$
\Pr[\hw{z\A}=h_1\mid \hw{z}=u]
=
\frac{\enum_{u,h_1}}{\binom{n}{u}}
\le \left(4e^2\frac{h_1}{n}\right)^{u/2}.
$
If $8e^2\varepsilon<1$, then $4e^2(h_1/n)\le 8e^2\varepsilon<1$ and the RHS is
decreasing in $u$. Hence, for all $u\ge u_{\min}(w)$,
$
\Pr[\hw{z\A}=h_1\mid \hw{z}=u]
\le \left(4e^2\frac{h_1}{n}\right)^{u_{\min}(w)/2}.
$
Therefore, using $\sum_u \Pr[u\mid w]=1$ and the fact that $u\ge u_{\min}(w)$,
we may upper bound the $u$-average by the worst case:
$
\sum_{u}\Pr[u\mid w]\Pr[h_1\mid u]
\le
\left(4e^2\frac{h_1}{n}\right)^{u_{\min}(w)/2}$ for $ (h_1\le 2\varepsilon n).
$

Substituting into \eqref{eq:Elow_expand} and using \eqref{eq:stage_probs} for
the second accumulator yields
\begin{equation}
\label{eq:Elow_after_mono}
E_{\le \varepsilon n}
\le
\sum_{w=1}^{k}\binom{k}{w}
\sum_{h_1=0}^{n}\sum_{h=1}^{\lfloor \varepsilon n\rfloor}
\left(4e^2\frac{h_1}{n}\right)^{u_{\min}(w)/2}
\cdot
\frac{\enum_{h_1,h}}{\binom{n}{h_1}}.
\end{equation}


%------------------------------------------------------------
\noindent\textbf{Step 3: Truncation (support constraints).}

From \eqref{eq:acc_exact}, $\enum_{a,b}=0$ unless $\lceil a/2\rceil\le b$.
Thus $\frac{\enum_{h_1,h}}{\binom{n}{h_1}}=0$ unless $h_1\le 2h$.
Since $h\le \varepsilon n$, we may restrict to
\begin{equation}
\label{eq:h1_range}
1\le h_1 \le 2\varepsilon n.
\end{equation}
Also, the first stage has $\enum_{u,h_1}=0$ unless $u\le 2h_1$; since
$u\ge u_{\min}(w)$, we must have $u_{\min}(w)\le 2h_1$, i.e.
$
d'\Big\lceil \frac{w}{\sigma}\Big\rceil \le 2h_1
\quad\Rightarrow\quad
w \le \sigma\left\lfloor\frac{2h_1}{d'}\right\rfloor.
$
So \eqref{eq:Elow_after_mono} becomes
\begin{equation}
\label{eq:Elow_trunc}
E_{\le \varepsilon n}
\le
\sum_{h_1=1}^{\lfloor 2\varepsilon n\rfloor}
\sum_{w=1}^{\sigma\lfloor 2h_1/d'\rfloor}\binom{k}{w}
\left(4e^2\frac{h_1}{n}\right)^{u_{\min}(w)/2}
\sum_{h=1}^{\lfloor \varepsilon n\rfloor}
\frac{\enum_{h_1,h}}{\binom{n}{h_1}}.
\end{equation}

%------------------------------------------------------------
\noindent\textbf{Step 4: Bound the second-stage tail $\sum_{h\le \varepsilon n}\Pr[h\mid h_1]$.}

By Lemma~\ref{lem:small_eps}(i) with $u\leftarrow h_1$,
\[
\frac{\enum_{h_1,h}}{\binom{n}{h_1}}
\le
\left(4e^2\frac{h}{n}\right)^{h_1/2}.
\]
Therefore,
\begin{equation}
\label{eq:second_tail}
\sum_{h=1}^{\lfloor \varepsilon n\rfloor}\frac{\enum_{h_1,h}}{\binom{n}{h_1}}
\le
\sum_{h=1}^{\lfloor \varepsilon n\rfloor}\left(4e^2\frac{h}{n}\right)^{h_1/2}
\le
(\varepsilon n)\left(4e^2\varepsilon\right)^{h_1/2},
\end{equation}
where the last step uses monotonicity in $h$ and $h\le \varepsilon n$.

Substitute \eqref{eq:second_tail} into \eqref{eq:Elow_trunc}:
\begin{equation}
\label{eq:Elow_mid}
E_{\le \varepsilon n}
\le
(\varepsilon n)
\sum_{h_1=1}^{\lfloor 2\varepsilon n\rfloor}
\left(4e^2\varepsilon\right)^{h_1/2}
\sum_{w=1}^{\sigma\lfloor 2h_1/d'\rfloor}\binom{k}{w}
\left(4e^2\frac{h_1}{n}\right)^{u_{\min}(w)/2}.
\end{equation}

%------------------------------------------------------------
\medskip
\noindent\textbf{Step 5: Re-indexing by active blocks.}

For an input $\xv\in\{0,1\}^k$, let $b=b(\xv)\in\{1,\dots,k/\sigma\}$ denote
the number of $\sigma$-bit input blocks on which $\xv$ is nonzero.
Since each active block contributes output weight at least $d'$, we have
\begin{equation}
\label{eq:umin_blocks}
u=\hw{\xv\B}\;\ge\; d' b .
\end{equation}

For fixed $b$, the number of inputs activating exactly $b$ blocks is at most
\begin{equation}
\label{eq:count_blocks}
N_b \;\le\; \binom{k/\sigma}{b}(2^\sigma-1)^b,
\end{equation}
since we choose the $b$ active blocks and then choose any nonzero
$\sigma$-bit pattern in each active block.


We now bound $E_{\le \varepsilon n}$ by summing over inputs grouped by the number
$b=b(\xv)$ of active $\sigma$-bit blocks (rather than by total Hamming weight).
From \eqref{eq:Elow_start}, equivalently
\[
E_{\le \varepsilon n}
=
\sum_{\xv\in\{0,1\}^k\setminus\{0\}}
\Pr\!\big[\hw{\xv\G}\le \varepsilon n\big].
\]
Fix such an input $\xv$ and let $b=b(\xv)$. By \eqref{eq:umin_blocks} we have
$u=\hw{\xv\B}\ge d'b$. Using the support constraint $\enum_{u,h_1}=0$ unless
$u\le 2h_1$ and applying \eqref{eq:second_tail}, we obtain
\[
\Pr\!\big[\hw{\xv\G}\le \varepsilon n\big]
\le
(\varepsilon n)\sum_{h_1=1}^{\lfloor 2\varepsilon n\rfloor}
\left(4e^2\varepsilon\right)^{h_1/2}\Pr[h_1\mid u].
\]
For each $h_1\le 2\varepsilon n$, Lemma~\ref{lem:small_eps}(i) gives
\[
\Pr[h_1\mid u]=\frac{\enum_{u,h_1}}{\binom{n}{u}}
\le
\left(4e^2\frac{h_1}{n}\right)^{u/2}
\le
\left(4e^2\frac{h_1}{n}\right)^{d'b/2},
\]
where the last inequality uses $u\ge d'b$ and $4e^2(h_1/n)\le 8e^2\varepsilon<1$.
If $d'b>2h_1$ then $\Pr[h_1\mid u]=0$, so only $b\le 2h_1/d'$ contribute.
Summing over all inputs with exactly $b$ active blocks and using \eqref{eq:count_blocks}
gives
\begin{align}
E_{\le \varepsilon n}
&\le
(\varepsilon n)
\sum_{h_1=1}^{\lfloor 2\varepsilon n\rfloor}
\left(4e^2\varepsilon\right)^{h_1/2}
\sum_{b=1}^{\lfloor 2h_1/d'\rfloor}
N_b\left(4e^2\frac{h_1}{n}\right)^{d'b/2}.
\label{eq:Elow_bh1}
\end{align}



\medskip
\noindent\textbf{(5b) Bounding the $h_1$-sum and geometric decay in $b$.}

We now swap the order of summation. Since $b\le 2h_1/d'$ if and only if
$h_1\ge d'b/2$, \eqref{eq:Elow_bh1} becomes
\begin{align}
E_{\le \varepsilon n}
&\le
(\varepsilon n)
\sum_{b=1}^{\lfloor 4\varepsilon n/d'\rfloor}
N_b
\sum_{h_1=\lceil d'b/2\rceil}^{\lfloor 2\varepsilon n\rfloor}
\left(4e^2\varepsilon\right)^{h_1/2}
\left(4e^2\frac{h_1}{n}\right)^{d'b/2}.
\label{eq:swap_b_h1}
\end{align}

Define
$
f(h_1)
:=
\left(4e^2\varepsilon\right)^{h_1/2}
\left(4e^2\frac{h_1}{n}\right)^{d'b/2}.
$
For $h_1\ge d'b/2$, we bound the ratio
$
\frac{f(h_1+1)}{f(h_1)}
=
(4e^2\varepsilon)^{1/2}
\left(1+\frac{1}{h_1}\right)^{d'b/2}
\le
(4e^2\varepsilon)^{1/2}
\exp\!\left(\frac{d'b}{2h_1}\right)
\le
(4e^2\varepsilon)^{1/2}e,
$
where the last inequality uses $h_1\ge d'b/2$.
Choosing $\varepsilon>0$ sufficiently small so that
$(4e^2\varepsilon)^{1/2}e<1$, the inner sum is geometrically decreasing,
and there exists a constant $C_0$ (independent of $n$ and $b$) such that
\begin{equation}
\label{eq:h1_tail_by_first}
\sum_{h_1=\lceil d'b/2\rceil}^{\lfloor 2\varepsilon n\rfloor} f(h_1)
\;\le\;
C_0\, f(\lceil d'b/2\rceil).
\end{equation}

Using $\lceil d'b/2\rceil\le d'b$ and $h_1/n\le d'b/n$ at the lower limit,
we obtain
$
f(\lceil d'b/2\rceil)
\le
\left(4e^2\varepsilon\right)^{d'b/4}
\left(4e^2\frac{d'b}{n}\right)^{d'b/2}.
$
Substituting into \eqref{eq:swap_b_h1} yields
\begin{equation}
\label{eq:Elow_after_h1}
E_{\le \varepsilon n}
\le
(\varepsilon n)\,C_0
\sum_{b=1}^{\lfloor 4\varepsilon n/d'\rfloor}
N_b
\left(4e^2\varepsilon\right)^{d'b/4}
\left(4e^2\frac{d'b}{n}\right)^{d'b/2}.
\end{equation}

Finally, using \eqref{eq:count_blocks} and
$\binom{k/\sigma}{b}\le (e k/(\sigma b))^b$ with $k=Rn$, we obtain
\[
E_{\le \varepsilon n}
\le
(\varepsilon n)\,C_0
\sum_{b=1}^{\lfloor 4\varepsilon n/d'\rfloor}
\left[
\frac{eR(2^\sigma-1)}{\sigma b}
\cdot
\left(4e^2\varepsilon\right)^{d'/4}
\cdot
\left(4e^2 d' b\right)^{d'/2}
\cdot
n^{\,1-d'/2}
\right]^b.
\]
We now analyze the $b$-sum.  Write the final bound as
\[
E_{\le \varepsilon n}
\le
(\varepsilon n)\,C_0
\sum_{b=1}^{\lfloor 4\varepsilon n/d'\rfloor}
\left[
K(\sigma,R,d',\varepsilon)\cdot b^{d'/2-1}\cdot n^{\,1-d'/2}
\right]^b,
\]
where
\[
K(\sigma,R,d',\varepsilon)
:=
\frac{eR(2^\sigma-1)}{\sigma}\cdot
\left(4e^2\varepsilon\right)^{d'/4}\cdot
\left(4e^2 d'\right)^{d'/2}
\]
absorbs all factors independent of $b$ and $n$.

\paragraph{The bottleneck is $b=1$.}
The $b=1$ term contributes
\[
E_{\le \varepsilon n}^{(b=1)}
\;\le\;
(\varepsilon n)\,C_0\cdot K(\sigma,R,d',\varepsilon)\cdot n^{\,1-d'/2}
=
O\!\left(n^{\,2-d'/2}\right).
\]
Therefore, to ensure $E_{\le \varepsilon n}\to 0$ as $n\to\infty$, we require
\[
2-\frac{d'}{2}<0
\qquad\Longleftrightarrow\qquad
d'>4.
\]

\paragraph{The terms with $b\ge 2$ are negligible.}
Let $\alpha:=d'/2-1>1$. The $b$th summand equals
\[
\left[K(\sigma,R,d',\varepsilon)\left(\frac{b}{n}\right)^{\alpha}\right]^b.
\]
For all sufficiently large $n$ (so that $\sqrt{n}\le 4\varepsilon n/d'$), split the range into
$2\le b\le \lfloor \sqrt{n}\rfloor$ and $b\ge \lceil \sqrt{n}\rceil$.

If $2\le b\le \sqrt{n}$ then $b/n\le n^{-1/2}$, so
\[
\left[K\left(\frac{b}{n}\right)^{\alpha}\right]^b
\le
\left(K\,n^{-\alpha/2}\right)^b,
\]
and hence, for all sufficiently large $n$ (so that $K\,n^{-\alpha/2}<1/2$),
\[
\sum_{b=2}^{\lfloor \sqrt{n}\rfloor}
\left[K\left(\frac{b}{n}\right)^{\alpha}\right]^b
\le
\sum_{b=2}^{\infty}\left(K\,n^{-\alpha/2}\right)^b
=
O\!\left(n^{-\alpha}\right).
\]

For $b\ge \sqrt{n}$, use the truncation $b\le \lfloor 4\varepsilon n/d'\rfloor$ to get
$b/n\le 4\varepsilon/d'$. Define
\[
\rho:=K(\sigma,R,d',\varepsilon)\left(\frac{4\varepsilon}{d'}\right)^{\alpha},
\]
and choose $\varepsilon$ sufficiently small so that $\rho<1$. Then for all $b\ge \sqrt{n}$,
\[
\left[K\left(\frac{b}{n}\right)^{\alpha}\right]^b \le \rho^b \le \rho^{\sqrt{n}},
\]
so
\[
\sum_{b=\lceil \sqrt{n}\rceil}^{\lfloor 4\varepsilon n/d'\rfloor}
\left[K\left(\frac{b}{n}\right)^{\alpha}\right]^b
\le
\frac{\rho^{\sqrt{n}}}{1-\rho}
=o(n^{-2}).
\]
Multiplying by the outer factor $(\varepsilon n)C_0$ shows that the total contribution of
$b\ge 2$ is $o(1)$. Together with the $b=1$ estimate above, this implies
$E_{\le \varepsilon n}\to 0$ as $n\to\infty$.

\medskip
Finally, by Markov's inequality,
\[
\Pr[d_{\min}(\G)\le \varepsilon n]\le E_{\le \varepsilon n}\xrightarrow[n\to\infty]{}0,
\]
which proves the theorem.


\end{proof}



